\documentclass[12pt]{article}
\usepackage[utf8]{inputenc}
\usepackage[T1]{fontenc}
\usepackage{hyperref}
\title{Kayap\'o Language Resources for K-12 Curriculum Development}
\author{Compiled by Claude Sonnet 5 (claude-sonnet-5) for Andr\'as Kornai}
\date{}
\begin{document}
\maketitle
\section{Introduction}
Kayap\'o (Mebêngôkre, ISO 639-3 txu) is a Je language of Brazil. Searches found mainly Brazilian academic work on school materials. The search did not surface a published general-purpose school dictionary or textbook series.
\section{Primary Educational Resources}
\begin{itemize}
\item UFRJ master's dissertation presenting an illustrated bilingual Kayap\'o--Portuguese school glossary of local fauna, made with students in M\^oxk\`ark\^o village: \url{https://pantheon.ufrj.br/handle/11422/29106}
\end{itemize}
\section{Linguistic and Reference Materials}
\begin{itemize}
\item Dictionary prototype of the B\~ep naming-ceremony lexicon, Mebêngôkre--Portuguese (UFRJ thesis, 2018): \url{https://bdtd.ibict.br/vufind/Record/UFRJ_f0c31687ab54484d9d14d120f287bc73}
\end{itemize}
\section{Religious and Community Resources}
None found in search results.
\section{Digital and Audio Resources}
None found in search results.
\section{Curriculum Development Context}
In 2019 the language was made co-official in S\~ao F\'elix do Xingu (Par\'a), according to an Abralin report, which also notes that this has not yet changed the local linguistic reality: \url{https://cadernos.abralin.org/index.php/cadernos/article/download/687/845/12791}. Existing materials are prototype or single-village efforts, so a curriculum would build on them and on community teachers.
\end{document}